\scanpage{119}
\mspage{59}
celui d'un rev. universel \deleted{\unsure{(ou}} \inserted{ou encore} d'un pt géom $\unsure{\xi|\zeta}$ de $\unsure{\widetilde K_x|\widetilde U_x}$ --- i.e. \unsure{d'une|de une} clôture algébrique de $\widetilde K_x$ --- \uline{et} d'une \deleted{\unsure{isom}} ``chemin de \unsure{domaines|connexion}'' entre $\unsure{\xi|\zeta}_{\widetilde U_x}\unsure{\mathrel{\text{---}}|\longrightarrow}\unsure{\widetilde K_x|\widetilde U_x}\longrightarrow\operatorname{Spec}(L)$ et $\unsure{\xi|\zeta}_L$ défini par $\unsure{\widetilde L|\overline L}$ --- ou ce qui revient au \unsure{=|même}, \inserted{à isom près des \unsure{revts}} d'un $\unsure{K|L}$-hom de \deleted{$\unsure{K}$} $\widetilde L_x$ dans $\overline L$ --- et ceci à son tour équivaut à la donnée d'un $x'$, \deleted{\uline{et} \unsure{d'un is}} comme dessus. Ceci posé, on a des isom can.
\[
(30)\quad\left\{\begin{aligned}
N_{x'}&\simeq\mathcal E_{\widetilde U_x}^{\overline L}\\
I_{x'}&\simeq\text{Noyau de l'hom canonique }\inserted{\text{surjectif}}\\
&\mathcal E_{\widetilde U_x}^{\overline L}\longrightarrow\unsure{\mathcal E_{\overline{k(x)}}|\mathcal E_{k(x)}^{\overline{k(x)}}}
\end{aligned}\right.
\]
(où la clôture alg $\overline{k(x)}$ de $k(x)$ est induite par la clôture alg. $\overline L$ de $\widetilde K_x$)
\margin{Cf théorie galoisienne des traits henséliens. Ici $I_{x'}$ joue le rôle d'un $\pi_{\widetilde U_x/k(x)}^{\overline L}$ ; en fait il est $=\pi_{\widetilde U_x/\widetilde S_x}^{\overline L}$ !}
(\unsure{Alors} (27) est déduit par théorie de Kummer)
\deleted{$I_{x'}\simeq$}
\deleted{(\unsure{27|31})}

\margin{$\Vert$ ?}
Je présume que (dans la situation actuelle, ou même, avec un corps de base $K$) on a $N_{x'}=\operatorname{Norm}_{\mathcal E_L}(I_{x'})$, \inserted{(a priori, on a seulement $\subset$)} ce qui est lié à l'idée qu'on doit avoir $I_{x'}=\operatorname{Norm}_{\pi_{L/K}}(I_{x'})$
\margin{118}
